\chapter*{阅读约定}
本文从每轮外层的 Mirror Descent 子问题出发，依次定义冻结信号 $f$ 与 advantage，求出加权教师 $m$ 的显式表达式，再反推产生该镜像方向的教师目标 $J$。前三章完成这条推导；第四章介绍 inner 如何拟合已经确定的 $m$，第五、六章分析参数性质与有限拟合。附录对每项独立结论先陈述命题与条件，再给完整证明。

镜像下降（Mirror Descent，MD）在概率空间通过局部线性目标与 Bregman 正则定义更新。本文沿用来源中的 MDPO、BPO、PMD-Part、PMD-Mean 名称，依次表示镜像目标优化、中心化残差回归、完整最优性方程回归和均值近似回归；第~\ref{ch:inner} 章给出它们的精确定义。

精确分布更新、总体拟合目标、有限样本估计和神经网络参数更新是四个层次。定理在哪一层成立，就在哪一层表述。完整词表求和消除条件动作积分误差，仍保留前缀覆盖与模型表示误差。本文的解析结论不直接给出任务正确率排序。

\begin{center}
\begin{tabular}{lp{100mm}}
\toprule
记号 & 固定含义\\
\midrule
$T,q_k,p_{k,n}$ & 教师、第 $k$ 轮冻结锚点、第 $n$ 步内层学生\\
$m_k$ & 该轮由最优化问题定义的目标教师\\
$\kappa,\eta$ & 目标形状、教师插值权重；正文取 $0\le\kappa\le1,\ 0<\eta<1$\\
$\alpha_{\rm opt},\tau$ & 参数优化器学习率、固定 $J_{\kappa,\eta}$ 后另选的镜像步长\\
$t,t_{\rm in}$ & 内层梯度流时间、每轮固定内层运行时长\\
$F_k,J_{\kappa,\eta},L_k$ & 重心目标、教师目标、内层拟合损失\\
$f,A=f/\eta,K_T$ & 冻结信号、标准化教师信号、冻结教师中心\\
$\mathcal A_k=C_qA$ & 本轮零条件均值的 advantage；$g_k=-\mathcal A_k$\\
$C_qh=h-\E_qh$ & 按同一前缀的锚点分布中心化\\
$d_k,\widehat d_k$ & 总体、经验前缀访问测度\\
$K_{\rm head},M,N$ & 头部动作数、同前缀辅助动作数、内层步数\\
\bottomrule
\end{tabular}
\end{center}

所有对数均为自然对数。$\TV(P,Q)=\frac12\sum_a|P_a-Q_a|$；平方 Hellinger 距离固定为 $H^2(P,Q)=\sum_a(\sqrt{P_a}-\sqrt{Q_a})^2$，不带 $1/2$。向量方差表示协方差矩阵的迹。正文省略前缀条件时，所有分布仍条件于同一个前缀。

附录按数学依赖组织：A 给出测度和基础工具；B 完整识别教师目标；C 处理内层投影和优化；D 分析形状与矩；E 处理积分信息与估计；F 将中心估计接回拟合；G 分析外层误差与序列边界；H 比较 TOP-D 的目标教师。短阅读版与完整稿使用同一份正文源码。
